\documentclass{article}
\usepackage[utf8]{inputenc}
\usepackage{geometry}
\geometry{a4paper, margin=1in}
\usepackage{hyperref}
\usepackage{amsmath}

\title{ESC204: Comprehensive Air Blade Prototyping, Simulation \& Verification Protocol}
\author{Jinesis Group \and Praxis Bot (Justin)}
\date{\today}

\begin{document}
\maketitle

\section{Project Context \& Overview}
The system has evolved from a time-based DC motor cart to a precision linear actuator utilizing a Nema 17 stepper motor and A4988 driver. This track system will carry an air blade mechanism designed to clear debris (e.g., wet leaves) off a sloped roof. To satisfy the ESC204 rubrics, we must prove both the mechanical translation (the track) and the fluid dynamics (the air blade) through multi-fidelity verification: mathematical models, computational fluid dynamics (CFD), and physical hardware testing.

\section{Packing List for Tomorrow's Prototyping}
To ensure the session is productive and yields rubric-compliant data, bring the following items. 

\subsection{What to BRING:}
\begin{itemize}
    \item \textbf{Continuous Air Source:} Hair dryer (with a "cool" setting) or shop vacuum exhaust.
    \item \textbf{Nozzle Fabrication Materials:} Cardboard, scissors, and duct tape to build a converging slit nozzle. (Alternatively, a 3D printed nozzle from Myhal).
    \item \textbf{Test Surface (High Friction):} Heavy-grit sandpaper or a real asphalt shingle. \textit{(Do not use smooth cardboard or grass, as they fail to mimic roof fluid dynamics).}
    \item \textbf{Measurement Tools:} Digital kitchen scale (for force testing) and a ruler/measuring tape.
    \item \textbf{Test Debris:} Actual leaves from outside (to be soaked in water).
    \item \textbf{Visualization Tools:} Sewing thread or yarn (for tuft testing), and tape.
    \item \textbf{Electronics:} The assembled stepper motor track, A4988 driver, microcontroller, and 12V DC wall adapter.
\end{itemize}

\subsection{What NOT to Bring:}
\begin{itemize}
    \item \textbf{Ball Pump:} Do not use a ball pump. It provides low-volume, pulsed bursts of air (high PSI, low CFM). An air blade requires continuous, planar, high-volume flow to create a boundary layer.
\end{itemize}

\section{Digital Validation: Math \& Simulations}
Before physical testing, establishing a theoretical baseline neutralizes stakeholder rebuttals ("Will air even work?"). 

\subsection{Mathematical Calculations}
Calculate the drag force ($F_D$) required to overcome the static friction (stiction) of a wet leaf.
\begin{equation}
    F_D = \frac{1}{2} \rho v^2 C_D A
\end{equation}
Where $\rho$ is air density, $v$ is the air velocity from the nozzle, $C_D$ is the drag coefficient of the leaf, and $A$ is the cross-sectional area. Compare the required $F_D$ to the output force of the nozzle.

\subsection{Online CFD Visualization}
Use a free browser-based tool like \textbf{SimScale} or \textbf{SolidWorks Flow Simulation} (available on Myhal lab computers).
\begin{itemize}
    \item \textbf{Setup:} Create a simple 2D simulation of a narrow nozzle blowing air across a flat boundary at a $30^\circ$ angle.
    \item \textbf{Goal:} Generate a velocity vector heat-map showing the "Coanda effect" (the air attaching to the surface) and the shear stress applied to the boundary layer. This visual proves the theoretical validity of the air blade.
\end{itemize}

\section{Physical Testing Protocols}

\subsection{Test 1: Output Force Quantification}
Tape a piece of cardboard vertically to the digital kitchen scale. Hold the nozzle horizontally and blow air at the cardboard from 10\,cm, 20\,cm, and 30\,cm. Record the mass reading in grams and convert to Newtons ($F = m \times 0.0098$). This establishes the actual force output of your continuous air source.

\subsection{Test 2: Angle of Attack Optimization}
Press wet leaves onto your textured sloped surface (sandpaper/shingle). Test the nozzle at $15^\circ$, $30^\circ$, and $45^\circ$ angles. Record which angle clears the leaves fastest. This demonstrates design optimization to the grading TAs.

\subsection{Test 3: Fluid Flow Visualization (Tuft Test)}
Tape 1-inch strips of thread across the test surface. Turn on the air blade. The tufts should lay flat against the surface, physically proving the air is attaching to the roof and creating a shear layer, corroborating the CFD simulation.

\section{Code \& Track Integration}
Once the stationary air blade is optimized, it must be integrated with the microcontroller code.
\begin{itemize}
    \item \textbf{Logic Integration:} Ensure the microcontroller code effectively utilizes the STEP and DIR pins of the A4988. 
    \item \textbf{Dynamic Testing:} Mount the optimized air nozzle to the moving track carriage. Run the stepper motor to translate the nozzle across the 60\,cm track while the continuous air source is active.
    \item \textbf{Verification:} Verify that the combined forward motion of the track and the planar flow of the air blade can clear a line of wet leaves without the track stalling under the resistance of the air hose.
\end{itemize}

\end{document}
